\chapter{Conclusions and Recommendations}
\label{chap:conclusion}

The project established a traceable engineering workflow for parameterized marine components. Bracket sampling and optimization reduced analytical mass from 0.4181~kg to 0.1126~kg while retaining predicted FOS 2.60 and deflection 0.361~mm. Engineering definitions and validation workflows were also completed for a steel lifting padeye and aluminium stabilizer.

The 73.1\% reduction is a screening result, not a validated design saving. The analytical model omits local and cyclic failure modes, and Fusion structural results were not recorded. The lightest bracket must therefore not be fabricated or installed on the basis of this report alone.

\section{Recommendations}
\begin{enumerate}
\item Complete mesh-converged FEA for every baseline and at least three finalists.
\item Add bolt bearing, tear-out, weld-group, local plate, fatigue, buckling, and corrosion checks.
\item Proof-load fabricated bracket and padeye specimens and measure strain and displacement.
\item Derive stabilizer pressure from CFD, model testing, or an accepted hydrodynamic method.
\item Apply relevant classification-society and statutory requirements before vessel use.
\end{enumerate}

\section{Future Work}
Padeye development should include pin contact, net-section failure, weld fatigue, and proof loading. Stabilizer development should combine hydrodynamic pressure, structural FEA, buckling, fatigue, and root-bearing assessment. Propeller work should begin only after a defensible thrust, torque, efficiency, cavitation, and pressure model is available.
